\cotefile{156-9}
\scanpage{34}
\mspage{30}
% archive sheet number: 33
% margin note written sideways
\margin{\uline{8 juillet}}
% source lines 1--3: heading
\uline{Digression 3} Magmas \unsure{munis} de \unsure{strates} \inserted{dites} \uline{compactes} (\uline{Magmas \unsure{compactologisés|compactologiques}}).

\uline{Notion de figure localement finie}

% source lines 4--7
% slanted
\margin{\unsure{La donnée} de $M_c$ s'appelle une ``\uline{compactologie}''\unsure{.}}
On se donne une partie $M_c$ de $M$, (structure supplémentaire) formée des \uline{strates} dites \deleted{``admissibles''} ``\uline{compactes}''.

\deleted{Ax} L'axiome \unsure{préalable} est le suivant

% source lines 8--10
\uline{Mag cpt 1}\quad \deleted{a) Si $X\in M$, \unsure{et} $\Rightarrow$ alors $\widetilde X$ fini.}
\[
\begin{aligned}
\text{a)}\quad &X\in M_c\ \Longrightarrow\ \widetilde X\ \text{fini.}\\
\text{b)}\quad &X\in M_c,\ Y\ll X\ \Longrightarrow\ Y\in M_c
\end{aligned}
\]
% source lines 10--13
--- en d'autres termes : $M_c$ est fermé dans $(M,\ll)$, ce qui revient au même, dans $(M,\overset{\circ}{\ll})$ et dans $(M,\trianglelefteq)$.

% slanted; left margin opposite lines 11--19
\margin{Un \unsure{magma|magasin} \unsure{compactologisé|compactologique} est dit ``à strates compactes'' si $M_c=M$. Il est dit \uline{compact} si $\exists F\in\unsure{\mathcal F|\mathscr F|\mathfrak F}(M)$, \unsure{à un} $\widetilde F$ \uline{fini}, tel que $\operatorname{supp}^{\circ}F=\unsure{M_c|M}$\unsure{.} i.e. $\nexists X\in M,\ F\mathrel{\mid\circ\mid}\deleted{\unsure{X}}X$.}
% slanted; left margin below the preceding note
\margin{\uline{Figure \unsure{compacte}} : $\widetilde F$ fini, et $\subset M_c$. \uline{\unsure{Préfigures|Préfigure}} ``\uline{compactes}'' : \deleted{\unsure{finies d'ordres}} une figure compacte.}

% source lines 14--17
Par la suite, on sera amené \unsure{sûrement} d'introduire un axiome Mag cpt 2, du type Borel-Lebesgue ; \inserted{et aussi que $M_c$ est ``riche'' dans $M$ (\unsure{p.v.|p.r.} aux subdivisions)}

% source lines 18--20
Quand on se donne de plus un \unsure{atelier} $\mathcal F$, \unsure{déterminant} un magasin, on exige

\deleted{Mag cpt (at)}

% source lines 21--24
\uline{Mag at cpt}\quad Si $F\ll X$, avec $F\in\mathcal F$, $X\in M_c$, alors $\widetilde F$ fini. Mieux : Si $X\in M_c$, alors $\forall F\in\mathcal F$, l'ensemble des $\deleted{X}Y\in\widetilde F$ telles que $\deleted{\widetilde F\mathrel{\mid\circ\mid}X}\ Y\parallel X$ est fini,
% left margin; continuation of the preceding sentence
\margin{i.e. $\widetilde F\mathrel{\mid\circ\mid}\widetilde X$, a un complémentaire fini dans $\widetilde F$.}

% source lines 25--28: paragraph cancelled with diagonal strokes
\deleted{On note ceci : Pour $M_c$ donné, \inserted{une compactologie donnée} il existe un $\mathcal F$ \unsure{attenant} à $M$, le plus grand possible \unsure{formé} \unsure{aux|avec|des} satisfaisant la condition envisagée. Les figures que}
